\documentclass{nrc2}
\usepackage{graphicx}
\usepackage{booktabs}
\usepackage{natbib}

\begin{document}

\title{Title of the Paper}

\author{First Author and Second Author}

\address{F. Author. Department of Earth Sciences, University, City, Country.\\
S. Author. Institute, City, Country.}

\correspond{F. Author (e-mail: first@example.org).}

\begin{abstract}
State the problem, what you did, and the main result in a few sentences.
\end{abstract}

\keywords{first keyword, second keyword, third keyword}

\maketitle

\section{Introduction}
\label{sec:intro}
Introduce the problem and why it matters, then the contributions of this
paper. Cite in parentheses~\citep{lorenz1963}, or by name, as \citet{manabe1967} showed, or several works at
once~\citep{lorenz1963,manabe1967}.

\section{Method}
\label{sec:method}
Describe what you did in enough detail to reproduce it. Numbered equations use
the \texttt{equation} environment:
\begin{equation}
  \label{eq:main}
  \frac{dx}{dt} = \sigma (y - x).
\end{equation}
Refer to it as Equation~(\ref{eq:main}).

\section{Results}
\label{sec:results}
Table~\ref{tab:results} and Figure~\ref{fig:overview} show the main results.

\begin{table}[htbp]
  \centering
  \caption{Results on the test set.}
  \label{tab:results}
  \begin{tabular}{lrr}
    \toprule
    Method & Accuracy (\%) & Time (s) \\
    \midrule
    Baseline & 81.2 & 12 \\
    Ours & \textbf{87.5} & 15 \\
    \bottomrule
  \end{tabular}
\end{table}

\begin{figure}[htbp]
  \centering
  \includegraphics[width=0.6\linewidth]{example-image}
  \caption{Overview of the method. Upload an image and use its file name here.}
  \label{fig:overview}
\end{figure}

\section{Discussion}
Interpret the results, compare them with earlier work and state the limitations.

\section{Conclusion}
Summarise the findings and what comes next.

\section*{Acknowledgements}
Thank the people and funding that supported this work.

\bibliographystyle{plainnat}
\bibliography{references}

\end{document}
